\chapter{Materials and Methods}\label{chapter:methods}
\thispagestyle{myheadings}

	% ========================================================================
	% MATERIALS AND METHODS
	%
	% This is the chapter that a biology / BME dissertation needs and that the
	% GRS Word template does not model. Neither the Mugar formatting guide nor
	% the GRS template prescribes any chapter structure -- that comes from your
	% committee -- but the skeleton below follows the standard order used in
	% neurophotonics and systems-neuroscience dissertations.
	%
	% Two rules of thumb worth respecting:
	%
	%  1. WRITE FOR REPLICATION. Every number a person would need in order to
	%     repeat the experiment belongs here: concentrations, volumes,
	%     coordinates, wavelengths, powers, frame rates, durations, ages,
	%     weights, catalogue numbers, software versions. If you find yourself
	%     writing "standard procedures were followed", replace it with the
	%     procedure.
	%
	%  2. MATCH THE MANUSCRIPT, THEN EXPAND. A dissertation methods chapter is
	%     normally the union of the methods sections of your papers, restored
	%     to full length. Journals cut methods for space; here you have none of
	%     that pressure, so put back the pilot experiments, the failed
	%     variants, the calibration procedures and the decision rules.
	%
	% Use \SI{} for every quantity so that units are typeset consistently:
	%     \SI{2}{\micro\liter}   \SI{920}{\nano\meter}   \SI{30}{\hertz}
	%     \SIrange{8}{12}{\week} \SI{0.5}{\milli\gram\per\kilogram}
	%     \ang{45}               \SI{37}{\celsius}
	%
	% Delete every section below that does not apply to your work, and add
	% sections for techniques that are missing.
	% ========================================================================

	This chapter describes the animals, surgical preparations, instrumentation,
	experimental designs, and analysis procedures common to
	\cref{chapter:study-one,chapter:study-two}.
	Procedures specific to a single study are described in the corresponding
	chapter.

	% ------------------------------------------------------------------------
	\section{Animals and ethical approval}\label{section:methods:animals}

		TODO: State the approving body and protocol number, the guidelines
		followed, the species, strain and source, the genotypes, the sex
		distribution, the age and weight range at the time of experiment, and
		the housing conditions.

		A model paragraph, for you to adapt:

		\begin{quote}
			All procedures were approved by the \acrshort{iacuc} of Boston
			University (protocol TODO) and conducted in accordance with the
			\acrshort{nih} \emph{Guide for the Care and Use of Laboratory
			Animals}. Reporting follows the \acrshort{arrive} guidelines.
			Experiments used TODO adult mice of both sexes
			(\nanimals{\textrm{TODO}} animals; TODO female, TODO male), aged
			\SIrange{8}{16}{\week} and weighing
			\SIrange{20}{30}{\gram} at the time of surgery. Animals were
			housed TODO per cage on a TODO light/dark cycle at
			\SI{22}{\celsius} with food and water available \emph{ad libitum}.
		\end{quote}

		\subsection{Mouse lines}\label{section:methods:animals:lines}

			TODO: Give the full formal strain name, the common name, the JAX or
			MMRRC stock number, and the breeding scheme for each line. State
			how genotype was confirmed.

		\subsection{Randomization, blinding and exclusion criteria}

			TODO: State how animals were assigned to groups, whether the
			experimenter or analyst was blinded, and the criteria by which
			animals or sessions were excluded. Report how many were excluded
			and why. Reviewers and committees ask about this, and deciding it
			after the fact is much harder than writing it down now.

	% ------------------------------------------------------------------------
	\section{Surgical procedures}\label{section:methods:surgery}

		\subsection{Anesthesia and analgesia}

			TODO: Induction and maintenance agents with doses and routes,
			depth-of-anesthesia monitoring, body temperature maintenance, eye
			protection, and the peri- and post-operative analgesia regimen
			with doses and durations.

		\subsection{Cranial window implantation}

			TODO: Stereotaxic coordinates relative to bregma, craniotomy
			diameter and location, dura handling, window material and
			thickness, sealing method, headpost or headplate type, and the
			recovery period before imaging.

		\subsection{Viral injections}\label{section:methods:surgery:virus}

			TODO: For each construct give the full name, the serotype, the
			source and catalogue or Addgene number, the titre in genome copies
			per millilitre, the injection coordinates and depths, the volume,
			the infusion rate, the pipette type, the wait time before
			withdrawal, and the expression interval before experiments.

			% Example table. Delete if unused -- but if you keep ANY table,
			% the List of Tables in main.tex is required.
			\begin{table}[htbp]
				\centering
				\caption[Viral constructs]{
					Viral constructs used in this dissertation.
					TODO: replace with your own.
				}\label{table:methods:constructs}
				\begin{tabular}{llll}
					\toprule
					Construct & Serotype & Titre (GC/mL) & Source \\
					\midrule
					TODO & TODO & TODO & TODO \\
					TODO & TODO & TODO & TODO \\
					\bottomrule
				\end{tabular}
			\end{table}

		\subsection{Chronic preparation and habituation}

			TODO: Handling schedule, habituation to head fixation, number and
			duration of sessions before data collection, and the criteria an
			animal had to meet before being included.

	% ------------------------------------------------------------------------
	\section{Optical instrumentation}\label{section:methods:optics}

		\subsection{Two-photon microscope}

			TODO: Describe the microscope as a light path. Laser make, model
			and tuning range; excitation wavelength and pulse width; power
			measured at the sample and how it was measured; objective
			magnification, \acrshort{na} and immersion medium; scanner type;
			dichroics and emission filters with centre wavelengths and
			bandwidths; detector type; and the acquisition software with
			version.

			TODO: State the resulting resolution explicitly, e.g. lateral and
			axial \acrshort{psf} full width at half maximum measured on
			sub-resolution beads.

		\subsection{Wide-field and intrinsic optical imaging}

			TODO: Illumination source and wavelengths, camera model, sensor
			binning, exposure, frame rate, and the resulting
			\si{\micro\meter\per\pixel} calibration.

		\subsection{Optogenetic and chemogenetic stimulation}

			TODO: For optogenetics: wavelength, source, fibre or free-space
			delivery, spot size, irradiance at the tissue in
			\si{\milli\watt\per\square\milli\meter}, pulse width, frequency,
			train duration, and duty cycle. For chemogenetics: compound, dose,
			vehicle, route, and the interval between administration and
			recording, plus the vehicle-only control.

		\subsection{Physiological monitoring}

			TODO: Heart rate, respiration, arterial oxygen saturation, body
			temperature and blood gases as applicable, with instruments and
			sampling rates. State the physiological ranges you required for a
			session to be included.

	% ------------------------------------------------------------------------
	\section{Experimental design}\label{section:methods:design}

		TODO: Describe a session from the animal's point of view: what
		happened, in what order, for how long, and how many times. Give trial
		counts, inter-trial intervals, block structure, stimulus parameters,
		and the order in which conditions were presented. State the control
		conditions and what they control for.

		TODO: State how the sample size was determined. If a power analysis was
		performed, give the assumed effect size, alpha, power and the resulting
		n. If the sample size followed prior work or practical constraints, say
		so honestly -- that is a normal and acceptable answer.

	% ------------------------------------------------------------------------
	\section{Data acquisition}\label{section:methods:acquisition}

		TODO: Sampling rates and bit depths for every channel, the
		synchronisation scheme between imaging, stimulation and behaviour,
		measured latencies or jitter, file formats, and total data volume.

	% ------------------------------------------------------------------------
	\section{Data processing}\label{section:methods:processing}

		TODO: Present the pipeline as an ordered list of transformations from
		raw files to the numbers that enter your statistics. Name every piece
		of software with its version, and state every parameter you did not
		leave at its default.

		\subsection{Motion correction and registration}

			TODO: Algorithm, reference image construction, rigid versus
			non-rigid, and the residual-motion criterion for discarding a
			session.

		\subsection{Segmentation and signal extraction}

			TODO: How \acrshortpl{roi} were defined -- manually, semi-automatically
			or by algorithm -- with the software and parameters. Neuropil or
			background subtraction and its coefficient. Bleaching correction.
			The exact definition of your normalised signal, written as an
			equation so there is no ambiguity:

			\begin{equation}\label{equation:methods:dff}
				\frac{\Delta F}{F_0}(t) = \frac{F(t) - F_0}{F_0},
			\end{equation}

			where $F(t)$ is the TODO-corrected fluorescence and $F_0$ is
			defined as TODO over the baseline window TODO.

		\subsection{Event detection}

			TODO: Detection method, threshold and how it was chosen, minimum
			event duration and amplitude, and how overlapping events were
			handled. Report the false-positive rate if you estimated one.

		\subsection{Quantified outcome measures}

			TODO: Define each dependent variable operationally. "Response
			amplitude" is not a definition; "mean $\Delta F/F_0$ over
			\SIrange{0.5}{2.0}{\second} after stimulus onset, minus the mean
			over the \SI{1}{\second} preceding it" is.

	% ------------------------------------------------------------------------
	\section{Statistical analysis}\label{section:methods:statistics}

		TODO: For each comparison state the test, the software, what the unit
		of analysis was, and how you handled the fact that cells, trials and
		sessions are nested within animals. Pooling cells across animals as if
		they were independent observations is the single most common
		statistical criticism in this field; if you used a mixed-effects model
		with animal as a random effect, say so here.

		TODO: State how normality and equal variance were assessed and what you
		did when they failed. Give the multiple-comparison correction and the
		family over which it was applied. State the alpha level. Say whether
		tests were one- or two-tailed.

		TODO: State what your error bars represent -- \acrshort{sem},
		\acrshort{sd} or \acrshort{ci} -- and use the same convention in every
		figure. Report exact $p$ values rather than inequalities, effect sizes
		with confidence intervals, and the $n$ for both animals and cells at
		every comparison.

	% ------------------------------------------------------------------------
	\section{Histology}\label{section:methods:histology}

		TODO: Perfusion protocol, fixative, post-fixation, sectioning method
		and thickness, antibodies with host species, catalogue numbers,
		dilutions and incubation conditions, counterstains, mounting medium,
		and the microscope used for verification. State how expression or
		targeting was confirmed and how many animals were excluded because it
		was not.

	% ------------------------------------------------------------------------
	\section{Data and code availability}\label{section:methods:availability}

		TODO: Repository, accession or DOI, licence, and what a reader would
		have to do to reproduce the figures. Many journals now require this,
		and committees increasingly ask for it.
